\documentclass[11pt,a4paper]{article}
\usepackage[utf8]{inputenc}
\usepackage[italian]{babel}
\usepackage{amsmath,amssymb,amsthm}
\usepackage{geometry}
\usepackage{enumitem}
\geometry{margin=2.3cm}

\newcommand{\R}{\mathbb{R}}
\newcommand{\dd}{\mathrm{d}}

\title{\textbf{Pool esame — Tier 1}\\\large 7 argomenti $\times$ 7 esercizi}
\author{Materiale didattico — uso interno}
\date{}

\begin{document}
\maketitle

\section*{07A — Formula di Taylor}

\textbf{1.} Scrivi il polinomio di Taylor di ordine 4 di $f(x)=\sin x$ centrato in $x_0=0$.\\
\emph{Sol.} $T_4(x)=x-\dfrac{x^3}{6}$ (il termine di grado 4 è nullo).

\textbf{2.} Sviluppa $f(x)=\ln(1+x)$ in $x_0=0$ fino al 5° ordine con resto di Peano.\\
\emph{Sol.} $\ln(1+x)=x-\dfrac{x^2}{2}+\dfrac{x^3}{3}-\dfrac{x^4}{4}+\dfrac{x^5}{5}+o(x^5)$.

\textbf{3.} Calcola $\displaystyle\lim_{x\to 0}\frac{e^x-1-x}{x^2}$ con Taylor.\\
\emph{Sol.} $e^x=1+x+\tfrac{x^2}{2}+o(x^2)\Rightarrow$ numeratore $=\tfrac{x^2}{2}+o(x^2)$, limite $=\tfrac{1}{2}$.

\textbf{4.} Calcola $\displaystyle\lim_{x\to 0}\frac{\sin x - x + \frac{x^3}{6}}{x^5}$.\\
\emph{Sol.} $\sin x = x-\tfrac{x^3}{6}+\tfrac{x^5}{120}+o(x^5)\Rightarrow$ numeratore $=\tfrac{x^5}{120}+o(x^5)$, limite $=\tfrac{1}{120}$.

\textbf{5.} Sviluppa $f(x)=\dfrac{1}{1-x}$ in $x_0=0$ fino al 4° ordine.\\
\emph{Sol.} Serie geometrica: $1+x+x^2+x^3+x^4+o(x^4)$.

\textbf{6.} Stima $\sqrt{1{,}08}$ con Taylor di $f(x)=\sqrt{1+x}$ in $x_0=0$ (ordine 2).\\
\emph{Sol.} $T_2(x)=1+\tfrac{x}{2}-\tfrac{x^2}{8}\Rightarrow T_2(0{,}08)=1+0{,}04-0{,}0008=1{,}0392$ (esatto $\approx 1{,}0392$).

\textbf{7.} Calcola $\displaystyle\lim_{x\to 0}\frac{\cos x - 1 + \frac{x^2}{2}}{x^4}$.\\
\emph{Sol.} $\cos x = 1-\tfrac{x^2}{2}+\tfrac{x^4}{24}+o(x^4)\Rightarrow$ numeratore $=\tfrac{x^4}{24}+o(x^4)$, limite $=\tfrac{1}{24}$.

\medskip\hrule\medskip

\section*{08 — Studio di funzione}

\textbf{1.} Studia $f(x)=\dfrac{x^2-1}{x^2+1}$.\\
\emph{Sol.} $D=\R$, pari, asintoto orizzontale $y=1$, $f'(x)=\tfrac{4x}{(x^2+1)^2}$: minimo in $x=0$, $f(0)=-1$. Immagine $[-1,1)$.

\textbf{2.} Studia $f(x)=xe^{-x}$.\\
\emph{Sol.} $D=\R$, $\lim_{x\to+\infty}f=0^+$, $\lim_{x\to-\infty}f=-\infty$, $f'(x)=(1-x)e^{-x}$: max in $x=1$, $f(1)=1/e$. Flesso in $x=2$.

\textbf{3.} Studia $f(x)=\dfrac{x^3}{x^2-1}$ (asintoto obliquo).\\
\emph{Sol.} $D=\R\setminus\{\pm 1\}$, dispari, asintoti verticali $x=\pm 1$, asintoto obliquo $y=x$. $f'(x)=\tfrac{x^2(x^2-3)}{(x^2-1)^2}$: max locale in $x=-\sqrt 3$, min locale in $x=\sqrt 3$.

\textbf{4.} Studia $f(x)=\ln(x^2+1)$.\\
\emph{Sol.} $D=\R$, pari, $f(0)=0$ minimo assoluto, $\lim_{x\to\pm\infty}=+\infty$. $f''(x)=\tfrac{2(1-x^2)}{(x^2+1)^2}$: flessi in $x=\pm 1$.

\textbf{5.} Studia $f(x)=\dfrac{x^2}{x-1}$.\\
\emph{Sol.} $D=\R\setminus\{1\}$, asintoto verticale $x=1$, asintoto obliquo $y=x+1$. $f'(x)=\tfrac{x(x-2)}{(x-1)^2}$: max locale in $x=0$ ($f=0$), min locale in $x=2$ ($f=4$).

\textbf{6.} Studia $f(x)=x^2\ln x$ su $(0,+\infty)$.\\
\emph{Sol.} $\lim_{x\to 0^+}f=0$, $f(1)=0$, $\lim_{x\to+\infty}=+\infty$. $f'(x)=x(2\ln x+1)$: minimo in $x=e^{-1/2}$, $f=-\tfrac{1}{2e}$.

\textbf{7.} Studia $f(x)=\dfrac{e^x}{x}$.\\
\emph{Sol.} $D=\R\setminus\{0\}$, asintoto verticale $x=0$ (a sinistra $-\infty$, a destra $+\infty$), asintoto orizzontale $y=0$ a $-\infty$. $f'(x)=\tfrac{e^x(x-1)}{x^2}$: minimo in $x=1$, $f(1)=e$.

\medskip\hrule\medskip

\section*{11 — Integrali doppi}

\textbf{1.} Calcola $\displaystyle\iint_D xy\,\dd A$ su $D=[0,1]\times[0,2]$.\\
\emph{Sol.} $\int_0^1 x\,\dd x \int_0^2 y\,\dd y = \tfrac{1}{2}\cdot 2 = 1$.

\textbf{2.} $\displaystyle\iint_D (x+y)\,\dd A$ sul triangolo di vertici $(0,0),(1,0),(0,1)$.\\
\emph{Sol.} $D:\,0\le y\le 1-x,\,0\le x\le 1$. $\int_0^1\!\int_0^{1-x}(x+y)\,\dd y\,\dd x=\int_0^1\Big[x(1-x)+\tfrac{(1-x)^2}{2}\Big]\dd x=\tfrac{1}{3}$.

\textbf{3.} $\displaystyle\iint_D \frac{1}{1+x^2+y^2}\,\dd A$ sul disco $x^2+y^2\le 1$.\\
\emph{Sol.} Polari: $\int_0^{2\pi}\!\int_0^1 \tfrac{\rho}{1+\rho^2}\dd\rho\,\dd\theta=2\pi\cdot\tfrac{1}{2}\ln 2=\pi\ln 2$.

\textbf{4.} Volume sotto $z=4-x^2-y^2$ sopra il disco $x^2+y^2\le 4$.\\
\emph{Sol.} Polari: $\int_0^{2\pi}\!\int_0^2 (4-\rho^2)\rho\,\dd\rho\,\dd\theta = 2\pi\big[2\rho^2-\tfrac{\rho^4}{4}\big]_0^2 = 2\pi\cdot 4 = 8\pi$.

\textbf{5.} Inverti l'ordine e calcola $\displaystyle\int_0^1\!\int_y^1 e^{x^2}\,\dd x\,\dd y$.\\
\emph{Sol.} Dominio: $0\le y\le x\le 1$. Invertito: $\int_0^1\!\int_0^x e^{x^2}\,\dd y\,\dd x = \int_0^1 x e^{x^2}\,\dd x = \tfrac{e-1}{2}$.

\textbf{6.} $\displaystyle\iint_D y\,\dd A$ con $D=\{(x,y):0\le x\le 1,\,x^2\le y\le x\}$.\\
\emph{Sol.} $\int_0^1\!\int_{x^2}^x y\,\dd y\,\dd x=\int_0^1 \tfrac{x^2-x^4}{2}\dd x=\tfrac{1}{2}\big(\tfrac{1}{3}-\tfrac{1}{5}\big)=\tfrac{1}{15}$.

\textbf{7.} Area della regione $D$ delimitata da $y=x^2$ e $y=2x$.\\
\emph{Sol.} Intersezioni $x=0,2$. Area $=\int_0^2 (2x-x^2)\dd x=\big[x^2-\tfrac{x^3}{3}\big]_0^2=\tfrac{4}{3}$.

\medskip\hrule\medskip

\section*{12 — Equazioni differenziali}

\textbf{1.} Risolvi $y'=xy$, $y(0)=1$.\\
\emph{Sol.} Variabili separabili: $\ln|y|=\tfrac{x^2}{2}+C\Rightarrow y=Ke^{x^2/2}$. Cauchy: $K=1$, $y=e^{x^2/2}$.

\textbf{2.} $y'+2y=e^x$.\\
\emph{Sol.} Lineare 1° ordine, $\mu=e^{2x}$: $(ye^{2x})'=e^{3x}\Rightarrow ye^{2x}=\tfrac{e^{3x}}{3}+C\Rightarrow y=\tfrac{e^x}{3}+Ce^{-2x}$.

\textbf{3.} $y''-3y'+2y=0$.\\
\emph{Sol.} Eq. caratt. $\lambda^2-3\lambda+2=0\Rightarrow\lambda=1,2$. $y=C_1 e^x+C_2 e^{2x}$.

\textbf{4.} $y''+y=\cos x$ (risonanza).\\
\emph{Sol.} Omogenea: $y_h=C_1\cos x+C_2\sin x$. Particolare: $y_p=\tfrac{x}{2}\sin x$. Generale: $y=C_1\cos x+C_2\sin x+\tfrac{x}{2}\sin x$.

\textbf{5.} Cauchy: $y'+\dfrac{y}{x}=x$, $y(1)=2$ ($x>0$).\\
\emph{Sol.} $\mu=x$: $(xy)'=x^2\Rightarrow xy=\tfrac{x^3}{3}+C$. Da $y(1)=2$: $C=\tfrac{5}{3}$, $y=\tfrac{x^2}{3}+\tfrac{5}{3x}$.

\textbf{6.} $y''-4y'+4y=0$ (radici doppie).\\
\emph{Sol.} $\lambda^2-4\lambda+4=0\Rightarrow\lambda=2$ doppia. $y=(C_1+C_2 x)e^{2x}$.

\textbf{7.} $y''+4y=\sin(3x)$.\\
\emph{Sol.} $y_h=C_1\cos(2x)+C_2\sin(2x)$. Particolare: $y_p=A\sin(3x)+B\cos(3x)$, sostituendo $A=-\tfrac{1}{5}$, $B=0$. $y=C_1\cos 2x+C_2\sin 2x-\tfrac{1}{5}\sin 3x$.

\medskip\hrule\medskip

\section*{13 — Funzioni di più variabili}

\textbf{1.} Determina il dominio di $f(x,y)=\dfrac{\sqrt{x-y^2}}{\ln(x+y)}$.\\
\emph{Sol.} $\{x\ge y^2\}\cap\{x+y>0\}\cap\{x+y\ne 1\}$.

\textbf{2.} Descrivi le curve di livello di $f(x,y)=x^2-4y^2$.\\
\emph{Sol.} $k=0$: $y=\pm x/2$ (rette). $k>0$: iperboli con asse trasverso su $x$. $k<0$: iperboli con asse su $y$.

\textbf{3.} $\displaystyle\lim_{(x,y)\to(0,0)}\dfrac{xy}{\sqrt{x^2+y^2}}$.\\
\emph{Sol.} Polari: $\rho\cos\theta\sin\theta$, $|\cdot|\le\rho\to 0$. Limite $=0$.

\textbf{4.} $\displaystyle\lim_{(x,y)\to(0,0)}\dfrac{x^2 y}{x^4+y^2}$.\\
\emph{Sol.} Lungo $y=mx$: $\to 0$. Lungo $y=x^2$: $\tfrac{x^4}{2x^4}=\tfrac{1}{2}$. Limite NON esiste.

\textbf{5.} La funzione $f(x,y)=\dfrac{x^3+y^3}{x^2+y^2}$ si estende con continuità in $(0,0)$?\\
\emph{Sol.} Polari: $\rho(\cos^3\theta+\sin^3\theta)$, $|\cdot|\le 2\rho\to 0$. Sì, $f(0,0)=0$.

\textbf{6.} Dominio di $f(x,y)=\arcsin(x+y)+\sqrt{4-x^2-y^2}$.\\
\emph{Sol.} $\{|x+y|\le 1\}\cap\{x^2+y^2\le 4\}$: striscia $-1\le x+y\le 1$ intersecata col disco di raggio 2.

\textbf{7.} Curve di livello di $f(x,y)=e^{x-y}$.\\
\emph{Sol.} $e^{x-y}=k>0\iff x-y=\ln k$: famiglia di rette parallele di pendenza 1.

\medskip\hrule\medskip

\section*{14 — Derivate parziali}

\textbf{1.} $f(x,y)=x^2\sin(xy)$. Calcola $f_x,f_y$.\\
\emph{Sol.} $f_x=2x\sin(xy)+x^2 y\cos(xy)$,\quad $f_y=x^3\cos(xy)$.

\textbf{2.} Gradiente di $f(x,y)=e^{x+y^2}$ in $(1,2)$.\\
\emph{Sol.} $\nabla f=(e^{x+y^2},2y\,e^{x+y^2})$, in $(1,2)$: $(e^5, 4e^5)$.

\textbf{3.} Derivata direzionale di $f(x,y)=x^2+xy+y^2$ in $(1,1)$ lungo $\mathbf u=(3,4)$.\\
\emph{Sol.} $\nabla f(1,1)=(3,3)$, versore $(\tfrac{3}{5},\tfrac{4}{5})$. $D_{\mathbf v}f=\tfrac{9+12}{5}=\tfrac{21}{5}$.

\textbf{4.} Piano tangente a $z=x^3+y^3$ in $(1,1,2)$.\\
\emph{Sol.} $f_x=3x^2,\,f_y=3y^2$, in $(1,1)$: $3,3$. $z=2+3(x-1)+3(y-1)=3x+3y-4$.

\textbf{5.} Verifica che $f(x,y)=\ln(x^2+y^2)$ soddisfa $f_{xx}+f_{yy}=0$ (Laplace).\\
\emph{Sol.} $f_{xx}=\tfrac{2(y^2-x^2)}{(x^2+y^2)^2}$, $f_{yy}=\tfrac{2(x^2-y^2)}{(x^2+y^2)^2}$, somma $=0$.

\textbf{6.} Stima $\sqrt{(3{,}02)^2+(4{,}01)^2}$ con linearizzazione di $f(x,y)=\sqrt{x^2+y^2}$ in $(3,4)$.\\
\emph{Sol.} $f(3,4)=5$, $f_x=\tfrac{3}{5}$, $f_y=\tfrac{4}{5}$. $f\approx 5+0{,}6\cdot 0{,}02+0{,}8\cdot 0{,}01=5{,}020$.

\textbf{7.} Verifica Schwarz per $f(x,y)=x^3 y^2 + \sin(xy)$.\\
\emph{Sol.} $f_{xy}=f_{yx}=6x^2 y+\cos(xy)-xy\sin(xy)$.

\medskip\hrule\medskip

\section*{15 — Massimi e minimi in più variabili}

\textbf{1.} Classifica i punti stazionari di $f(x,y)=x^3-12xy+8y^3$.\\
\emph{Sol.} $\nabla f=(3x^2-12y,-12x+24y^2)=0\Rightarrow$ punti $(0,0)$ e $(2,1)$. Hessiana: $\det H=12x\cdot 48 y-144$. $(0,0)$: sella. $(2,1)$: $\det>0$, $f_{xx}=12>0\Rightarrow$ minimo.

\textbf{2.} Estremi di $f(x,y)=x^2+xy+y^2-3x$.\\
\emph{Sol.} $\nabla f=(2x+y-3,x+2y)=0\Rightarrow x=2,y=-1$. $\det H=4-1=3>0$, $f_{xx}=2>0$: minimo, $f=-3$.

\textbf{3.} Max/min assoluti di $f(x,y)=x^2+y^2-x$ sul disco $x^2+y^2\le 1$.\\
\emph{Sol.} Interno: $\nabla f=0\Rightarrow (1/2,0)$, $f=-1/4$. Bordo: parametrizzo, $g(\theta)=1-\cos\theta$, max in $\theta=\pi$ ($f=2$), min in $\theta=0$ ($f=0$). Max assoluto $2$ in $(-1,0)$, min assoluto $-1/4$ in $(1/2,0)$.

\textbf{4.} Lagrange: max e min di $f(x,y)=x+2y$ sul vincolo $x^2+y^2=5$.\\
\emph{Sol.} $1=2\lambda x,\,2=2\lambda y\Rightarrow y=2x$. $x^2+4x^2=5\Rightarrow x=\pm 1$. $f(1,2)=5$, $f(-1,-2)=-5$.

\textbf{5.} Lagrange: estremi di $f(x,y)=xy$ sul vincolo $\tfrac{x^2}{4}+y^2=1$.\\
\emph{Sol.} $y=\lambda x/2$, $x=2\lambda y\Rightarrow x^2=4y^2$. Vincolo: $y^2+y^2=1\Rightarrow y=\pm\tfrac{1}{\sqrt 2}$, $x=\pm\sqrt 2$. Max $=1$ in $(\sqrt 2,\tfrac{1}{\sqrt 2})$ e $(-\sqrt 2,-\tfrac{1}{\sqrt 2})$; min $=-1$ negli altri due.

\textbf{6.} Max/min assoluti di $f(x,y)=(x-1)(y-1)$ su $K=[0,2]\times[0,2]$.\\
\emph{Sol.} Interno: $\nabla f=(y-1,x-1)=0\Rightarrow (1,1)$, $f=0$. Bordo: ai 4 lati il prodotto $(x-1)(y-1)$ raggiunge $\pm 1$ nei vertici. Max $=1$ in $(0,0),(2,2)$; min $=-1$ in $(0,2),(2,0)$.

\textbf{7.} Classifica i punti stazionari di $f(x,y)=x^4+y^4-4xy$.\\
\emph{Sol.} $4x^3-4y=0\Rightarrow y=x^3$. $4y^3-4x=0\Rightarrow x=y^3=x^9\Rightarrow x(x^8-1)=0$. Punti: $(0,0),(1,1),(-1,-1)$. $\det H=(12x^2)(12y^2)-16=144x^2y^2-16$. In $(0,0)$: $-16<0$, sella. In $(\pm 1,\pm 1)$: $144-16=128>0$, $f_{xx}=12>0$: minimi.

\end{document}
